\section{His Political Ideas}

\begin{quotex}
Ideas belong to no one; they are what they are. \flright{\textsc{Charles Maurras}, \textit{My Political Ideas}}

\end{quotex}
Tonight, I'm outlining the plan for translating a series of essays by Charles Maurras. This fulfills a promise and responds to a challenge. The promise was made in the comments section of Taki's Magazine\footnote{\url{http://takimag.com/\#axzz27L592q87}} a few years ago when it used to be a conservative resource instead of the journal of choice for hip urban racists that it is today. The challenge is from \textbf{Alain de Benoist} to compare and contrast the ideas of Evola and Maurras. It appears no one has the will or the talent to accomplish that.

That still does not fully answer the question of ``Why bother?" In his prime, Mauras was highly influential, not just in France, but also on the Iberian peninsula and across Latin America. Echoes of his thought still reverberate in Traditional Catholic circles through the SSPX\footnote{\url{http://www.sspx.org/}} and TFP\footnote{\url{http://www.tfp.org/}}, although they are far from the mainstream. Yet for the modern world, his ideas are not only irrelevant today, they are positively distasteful. Nevertheless, there may still be salvageable concepts if we leave aside the more contentious parts of his writings.

The French Republic prides itself on its secularism, but Maurras himself was a secularist. The difference is that Mauras includes the entire history of France in his vision, while the Republic looks is concerned with the moment. Like the liberals, he also claims to base his worldview on science. However, for the liberal, the individual is the atom, and family, culture, and nation are the results of his will which he may freely choose or reject. For Maurras, family, culture, and nation come first, they are the source of the individual's identity and hence demand his allegiance. For the liberal, there is no fixed human nature since he creates his own identity. Maurras, on the other hand, based his understanding of a fixed human nature on the philosophy of Thomas Aquinas.

There is no compromise possible here. To the modern mind, Maurras' ideas are necessarily restrictive, as they limit his freedom in the present, and they are reactionary since they tie him to the actual past whereas he wants to be liberated for the possible future. For them, this exercise can serve as a museum for discarded ideas, useful for pedagogy, but not much else. The real challenge, then, is for the so-called New Right. This will serve as a touchstone to evaluate its doctrines which hardly differ from liberalism considering its commitment to a functional polytheism and its tolerance for the aesthetically ugly and the unnatural. We hope there may still be those looking for ideas on the Right that derive from traditional sources rather than from foraging in the weeds of postmodern thought for imaginary gems as does the New Right.

Maurras was a poet and a member of the prestigious \textbf{Académie française} until his postwar conviction of collaborating with the enemy. Hence, his prose is written in beautifully expressive French so it is a shame to have to translate him as I will be unable to capture that.

Maurras was not often a systematic writer, so it is difficult to find a representative sample of his work despite his voluminous output. We considered two works. One was written late in his life, Order and Disorder, which fits into a theme we have been exploring recently. However, it is not broken up into sections so it is not so suitable for a blog. The other is My Political Ideas which does attempt to lay out his ideas somewhat systematically. Next Sunday, we will start with the chapter titled ``Principles" and then decide if it is worthwhile to continue. These can and should be compared to some of the metaphysical principles we have been writing about. Like Abraham, I am hoping for ten good men to be able to provide insightful and relevant commentary on these essays, although in reality I would be satisfied with one. The forthcoming principles are:

\begin{enumerate}
\item Truth 
\item Strength 
\item Order 
\item Authority 
\item Liberty 
\item Right and Law 
\item Property 
\item Heritage 
\item Duty of Heritage 
\item \ Tradition 
\end{enumerate}


\flrightit{Posted on 2012-09-23 by Cologero }

\begin{center}* * *\end{center}

\begin{footnotesize}\begin{sffamily}



\texttt{Matt on 2012-09-23 at 23:37 said: }

Excellent. The translation of Maurass' writing on Comte you provided a couple of years earlier was great so I'll be looking forward to this.

Both Maurras and Comte remain quite unique in their defense of traditional social-political structures while respectfully remaining non-committed to affirming the spiritual reality that the structures are intended to reflect. Though even then, Maurras became a committed Catholic near the end of his life, and Comte appeared to take a quasi-mystical outlook in his outlook with venerating the great predecessors that came before us.

Hope your health has improved.


\hfill

\texttt{Avery Morrow on 2012-09-24 at 00:52 said: }

``Tradition means giving votes to the most obscure of all classes, our ancestors. It is the democracy of the dead. " –Chesterton


\hfill

\texttt{Drieu on 2012-09-24 at 17:17 said: }

``However, for the liberal, the individual is the atom, and family, culture, and nation are the results of his will which he may freely choose or reject."

This is inaccurate and flat out stupid to anyone who has actually bothered to research the history of ideas. Not all liberals and leftists hold the same views, just like not all rightists hold the same views.

Immanuel Kant and Denis Diderot, both radicals in their time, did not see individuals as atoms but as *cultural agents*, i.e. that man was unavoidably shaped by his culture. The same apples to Johann Gottfried Herder, the philosophical founder of nationalism, who supported the French Revolution. 

The anarchist Mikhail Bakhunin wrote that ``Nationality is not a principle; it is a legitimate fact, just as individuality is. Every nationality, great or small, has the incontestable right to be itself, to live according to its own nature. This right is simply the corollary of the general principle of freedom."

Even among economic liberals, Milton Friedman explicitly states in Capitalism and Freedom that the basic component of society is not the individual, but the family.

Hell, Maurras's own ideas on nationalism (as well as Maurice Barrès's) were primarily shaped by Hippolyte Taine and Ernest Renan, both agnostic liberals (though of a conservative bent more akin to British Toryism). 

There are many more examples, but those are just some prominent ones off the top of my head. You need to drop the inane traditionalist platitudes and straw men about some monolithic ``Enlightenment" or ``liberalism" and actually figure out what the other side thinks. Especially since most of the ideas attributed to the ``Enlightenment" (secularism, republicanism, written constitutions, equality before the law, civil liberties) were organically developed well before then in the religious sphere by devout Christians, not by deists or skeptics.

``We hope there may still be those looking for ideas on the Right that derive from traditional sources rather than from foraging in the weeds of postmodern thought for imaginary gems as does the New Right."

Except that the right, including La Vera Destra to which your dear Evola refers, has always incorporated the latest philosophical trends into its thought. By absorbing some postmodern currents, the Nouvelle Droite has not at all betrayed the legacy of people like Joseph de Maistre or Juan Donoso Cortés. And considering that postmodernism thought aligns with certain motifs of the right, such as a suspicion of the Enlightenment narrative, disdain for rationalism, and rejection of an historical telos, I don't even see it as a necessarily corrupting element.

For the record, I'm not a leftist or liberal, but someone who is tired of seeing people who are nominally on my side repeat the same nonsense again and again and make themselves look like idiots in the process. 

As for Maurras, I'd go for one of the early works like Enquete sur la monarchie. You forgot to mention that his views changed a bit after WWI. In the early period, he advocated decentralization, regionalism, and flirted with unorthodox economic doctrines such a syndicalism. After the war, he shifted towards support for centralization while still retaining some affection for the provinces, and became more reconciled with capitalism and less interested in activism. The latter two changes ended up driving away many of his most talented disciples who were much more dedicated to the earlier radicalism, such as Georges Bernanos, Lucien Rebatet, Robert Brasillach, and Thierry Maulnier.


\hfill

\texttt{Andrew on 2012-09-24 at 18:10 said: }

And that is not the tradition we are talking about.


\hfill

\texttt{Cologero on 2012-09-24 at 19:03 said: }

Drieu, I'm afraid we differ in method. In order to understand ``liberalism" or any ``ism", it requires an understanding of the principles that govern it. Hence, we here reject out of hand the empirical approach of asking those who call themselves ``liberal" what their positions are on various issues. This fails for two reasons. The first is that focus turns to secondary issues. The second, and more important, is that few men are really aware of the principles and motives that drive their thought. In out language, we are interested in the ``thesis" rather than the ``hypothesis". Although I have hinted at that distinction previously, I will have to expand on it at a later time.

The sign of a great thinker is that he knows and can express his first principle. These are the men we need to focus on, not a hodge podge of random quotes. For example, Francis Bacon explicitly rejected ``formal causes", which is antithetical to any Traditional understanding of the world. This has consequences in what you call the ``history of ideas" even if later developments were not apparent to Bacon himself. Yet, this rejection is considered normal today even to the extent that very few men even realize that it underlies their entire worldview. We will deal with this again soon, probably on Wednesday.

Kant took this further by denying knowledge itself. Given Bacon's premise, Kant's conclusion absolutely follows. Yet this is opposed to any Traditional world view. So for Kant ``culture" is an empirical reality that ``shapes" an individual; unfortunately, that is not a very profound observation. The Traditional worldview is that a culture is a transcendental reality, something Kant denies necessarily as unknowable. It more than ``shapes" the individual; rather the individual is unthinkable without it.

So, yes, Drieu, we try to ``figure out" not just ``what" they think, but rather ``how" they think. And that is very often not exactly how they see themselves. In other words, what are the presuppositions of their thought? And, of course, we mean consistent thought since very few men are able to think consistently.

If this project has no appeal to you, you are welcome to leave and hang out wherever there are no stupid people or idiots. Traditional thought is always considered stupid, superstitious, and ignorant by the modern mind. Since you appear to have a modern mind, I don't see how we are on the ``same" side, not even nominally.

I am always amused by readers who want to tell me what to write on this blog, or what to translate, or what I omitted. Obviously, I omit quite a bit, since my contention is that a single blog post needs to be narrowly focused. If you prefer that I translate Enquete sur la monarchie for you, cut me a check for all the time required and I will move it to the very top of my priorities. Otherwise, I will select texts based on my own purposes and interests.


\hfill

\texttt{Jason-Adam on 2012-09-24 at 19:31 said: }

As someone whose life has been influenced by Maurras' ideas himself, I will gladly do what I can to assist in this project.

Let me ask Cologero, we know already the distinction between the traditional and modern minds. The new right seems to think the postmodern mind is something else and an improvement upon the modern that can lead back to traditional thinking. Do you agree and if not why ?


\hfill

\texttt{Izak on 2012-09-24 at 23:02 said: }

I'm not really seeing the willingness to discourse with postmodern theorists as any indication that the New Right is eo ipso postmodern, or informed by liberal values, or even definable in any meaningful way as of now.

The medieval Catholic justification for glossing heretics like Origen was that ``even in a heap of dung one can find a nugget of gold." 

Like most intellectual movements, the New Right will be defined by whatever influence it manages to concretely produce. Because that's what politics is all about.


\hfill

\texttt{Cologero on 2012-09-25 at 00:17 said: }

Jason-Adam, what I will be looking for are comments on the texts, either explicating the text, relating it to other ideas, or putting it into a broader context. I don't plan to ``debate" the New Right. I really do understand the allure of postmodernism; I have experienced it firsthand. The fundamental point is that if the knowledge of ultimate reality is already known, has been known since ancient times, and is handed down as tradition, then why the quest for novelty? 

I suggest to the advocates here of the new right that they explain why Benoist ultimately rejected Maurras … that would be useful and would add to the store of knowledge.

The other part of the project is to compare and contrast Maurras and Evola; comments in that regard are welcome. Moreover, I am curious why Benoist thinks it a worthwhile project as well as the reason neither he nor anyone else has followed through with it.


\hfill

\texttt{Avery on 2012-09-25 at 01:33 said: }

I employ vulgar sources wherever they may help people understand, and resort to fleshing out the Guenonian depths only when there has been a clear disruption of our common sense.

In this case, I see a clear connection between Chesterton's romantic definition and Cologero's statement, ``Maurras himself was a secularist. The difference is that Maurass includes the entire history of France in his vision, while the Republic looks is concerned with the moment."


\hfill

\texttt{Jason-Adam on 2012-09-25 at 20:25 said: }

I could write an essay explaining clearly the exact differences between Evola and Maurras so I will gladly comment when the texts appear. Let me just say one thing for now – that ultimately Maurras and de Benoist are looking to different pasts, one to the Middle Ages, the other to ancient Greece. I conjecture de Benoist has not gone through with the project because he would have to explicitly state then that he does admire mediaeval civilisation.


\hfill

\texttt{Jason-Adam on 2012-09-25 at 20:26 said: }

forgive me, it should say de Benoist does NOT admire mediaeval civilisation.


\hfill

\texttt{Dominion on 2012-09-26 at 00:50 said: }

The issue is one of form rather than essence, in my understanding. While it is true that the ultimate principles already exist, the way in which they manifest are not. This differs between Traditional civilizations; Christendom did not look identical to the Caliphate, which did not look identical to old Rome, and so on. Traditionalists in the New Right, particularly Dugin, use postmodern thought to critique the basic principles of Liberal civilization and thought, which is similar to what Maurras sought to do with positivist philosophy (``What our fathers did through custom and feeling, we ourselves pursue it through reason and will, with the assurance and clarity of science.") When Faye (who is admittedly not a Traditionalist) talks about reconciling Evola and Marinetti, I interpret this as taking the ultimate principles which we already know and seeing how we can live by them in our current, temporal and changing world. Tradition of principle must not be made equivilant to traditionalism of form. As Guido de Giorgio once said:

``Tradition is absolutely different from traditionalism: one is an eternally fecund living patrimony, rich with infinite potentialities in all times and circumstances . . . . the other is but sterile residue, an inefficient, self-enclosing concreteness impossible to adapt and lacking all energetic and creative force. Tradition is clearly opposed to traditionalism, just as truth is opposed to commonplaces."


\hfill

\texttt{Cologero on 2012-09-26 at 07:29 said: }

Well put, Dominion. The phrase you quoted also explains my own interest in the thought of Maurras and Comte. So are you saying that the Maurrasian critique of liberalism is no longer effective and that the postmodern critique may be more fruitful?


\hfill

\texttt{Dominion on 2012-09-26 at 14:46 said: }

Considering the influential role that positivism plays in the current order, particularly in the scientific sphere and even to a great extent in the social sciences, I think that both the Maurrasian and postmodern critiques of liberalism would prove useful. Personally, although I find the postmodern critiques interesting and useful, I am more inclined towards rational and scientific perspectives and might thus find Maurras more rewarding. I don't have much background in his thought, however, and thus await your translation.


\end{sffamily}\end{footnotesize}
